\chapter{État de l'art}
\label{ch:sota}

\section{Le rôle du Product Owner}
Défini dans le \textit{Scrum Guide}~\cite{schwaber2020scrum}, le PO est
responsable de la maximisation de la valeur produite par l'équipe via
la gestion ordonnée du Product Backlog. Ses compétences clés couvrent
la \emph{vision produit}, la \emph{priorisation}, la \emph{communication
transverse} et l'\emph{acceptation} des incréments. Il interagit avec
les parties prenantes pour collecter le besoin, avec l'équipe de
développement pour clarifier les stories, et avec le management pour
défendre les arbitrages.

\section{Frameworks Agile}
\begin{description}
  \item[Scrum] Rituels (Sprint Planning, Daily, Review, Rétrospective),
        artefacts (Product Backlog, Sprint Backlog, Incrément) et rôles
        (PO, Scrum Master, Dev Team).
  \item[\gls{safe}] Mise à l'échelle de Scrum sur plusieurs équipes via
        les Program Increments et le \gls{wsjf}~\cite{leffingwell2018safe}.
  \item[Kanban] Flux continu, limites de WIP (Work In Progress), mesure
        du lead time et du cycle time.
\end{description}

\section{Méthodes de priorisation}
\Cref{tab:prio-methods} compare les sept méthodes implémentées dans
POSTIE. Chacune répond à un contexte différent~: RICE pour les
arbitrages produit consumer, WSJF pour les contextes SAFe, MoSCoW
pour la communication aux parties prenantes, Kano pour l'analyse de
satisfaction, ICE pour les décisions rapides, Value/Effort pour la
visualisation matricielle, et ML Hybrid pour combiner les signaux.

\begin{table}[H]\centering\small
\caption{Comparaison des méthodes de priorisation implémentées.}
\label{tab:prio-methods}
\begin{tabularx}{\linewidth}{lXXl}
\toprule
\textbf{Méthode} & \textbf{Principe} & \textbf{Forces} & \textbf{Limites}\\
\midrule
\gls{rice} & (Reach\,$\times$\,Impact\,$\times$\,Confidence)/Effort
  & Quantitatif, populaire & Subjectivité du Reach\\
\gls{wsjf} & (BV+TC+RR)/JobSize
  & Adapté au flux de valeur & Estimation du Job Size\\
\gls{moscow} & Must / Should / Could / Won't
  & Simple, communicant & Faible granularité\\
Kano & Basique/Performance/Delighter
  & Centré utilisateur & Étude requise\\
\gls{ice} & Impact\,$\times$\,Confidence\,$\times$\,Ease
  & Rapide & Approximatif\\
Value/Effort & Matrice 2$\times$2 & Visuel & Binaire\\
ML~Hybrid & Blend pondéré des scores
  & Robuste & Calibration nécessaire\\
\bottomrule
\end{tabularx}
\end{table}

\section{RAG et LLMs appliqués au product management}
Introduite par Lewis~\textit{et al.}~\cite{lewis2020rag}, l'architecture
\gls{rag} consiste à augmenter un \gls{llm} par un système de
récupération vectorielle. Le flux est le suivant~: (1) la requête
utilisateur est encodée en vecteur, (2) le vector store retourne les
$k$ chunks de documents les plus pertinents, (3) ces chunks sont
injectés dans le prompt du LLM qui génère la réponse finale.

La stratégie \gls{mmr} permet de réduire la redondance des chunks
retournés en pénalisant les vecteurs trop similaires entre eux. Les
approches de routage multi-modèles exploitent les compromis
coût/latence/qualité entre différentes familles de LLMs~: un modèle
léger pour les tâches structurées, un modèle puissant pour le
raisonnement long.

\section{Positionnement de POSTIE}
POSTIE se distingue des outils existants (Jira Intelligence,
ProductBoard, Aha!) par~:
\begin{itemize}
  \item son caractère \emph{open-source} et auto-hébergeable~;
  \item la pluralité des algorithmes de priorisation accessibles dans
        une même interface~;
  \item le routage intelligent multi-modèles entre GPT-4o et
        Mistral Large~;
  \item la traçabilité documentaire systématique des réponses RAG~;
  \item l'intégration native d'un Sprint Planner sous contrainte de
        vélocité.
\end{itemize}
